%=====================================================================
\chapter{SEO and Digital Marketing}\label{ch:seo}
%=====================================================================
\locator{6}{Part VI --- Reach, Law and Practice}

\begin{outcomes}
By the end of this chapter you will be able to:
\begin{tight}
  \item \textbf{Distinguish} the parts of search optimisation you control and
        can check from the parts you cannot.
  \item \textbf{Measure} a shop's crawlability: URLs, titles, descriptions,
        canonical links, sitemap and first byte.
  \item \textbf{Explain} why readable URLs are two halves, and \textbf{apply}
        both.
  \item \textbf{Diagnose} a configuration change that makes a shop worse while
        appearing to improve it.
  \item \textbf{Prioritise} marketing work by what is checkable rather than by
        what is advised.
\end{tight}
\end{outcomes}

\begin{prereq}
\chref{security}: \texttt{AllowOverride None}, and why \texttt{.htaccess} is
never read on this stack. \chref{ux}: page weight and time to first byte.
\chref{catalogue}: the catalogue tables behind the pages.
\end{prereq}

\begin{keyterms}
crawler $\bullet$ indexing $\bullet$ readable URL $\bullet$ URL rewriting
$\bullet$ canonical link $\bullet$ meta description $\bullet$ title tag
$\bullet$ duplicate content $\bullet$ sitemap $\bullet$ robots.txt $\bullet$
faceted navigation $\bullet$ time to first byte $\bullet$ structured data
\end{keyterms}

\section{What Is Checkable, and What Is Not}

Most writing about search optimisation concerns things nobody outside the
search engine can verify: how a ranking is computed, how a change affected it,
what the algorithm rewards this month. This chapter does not go there, because
a course that cannot check its claims has no business making them.

\begin{conceptbox}[The question this chapter asks instead]
Not ``how do I rank higher'', but: \textbf{can a crawler read these pages, tell
them apart, and find all of them?}

\smallskip
Every part of that is checkable from outside, in about a minute, with the shop
in front of you. And a shop that fails it cannot rank at all --- so it is the
part worth doing first regardless of what you believe about the rest.
\end{conceptbox}

\section{Readable URLs Are Two Halves}

\begin{definitionbox}[Readable URL, URL rewriting]
A \term{readable URL} identifies a page by words rather than by parameters.
This course's demonstration category is reached at

\smallskip
\hspace*{1em}\texttt{/en-gb/catalog/desktops}

\smallskip
\noindent rather than at

\smallskip
\hspace*{1em}\texttt{index.php?route=product/category\&path=20}

\smallskip
It needs two independent things.

\smallskip
The \term{application half} generates them: the shop looks each link up in a
table of keywords and writes the readable form into its HTML.

\smallskip
The \term{server half} resolves them: the web server turns the readable path
back into the parameters the application understands, \emph{before} the
application sees the request.

\smallskip
Either half alone is useless. The second alone leaves the links unchanged. The
first alone is \emph{worse than nothing}.
\end{definitionbox}

\begin{measured}[Each half on its own, and both together]
From \texttt{code/m21\_seo.py --fix}. The shop ships with \textbf{78 prepared
keywords} in \texttt{oc\_seo\_url} and \texttt{config\_seo\_url} set to
\textbf{0}.

\rowcolors{2}{white}{lightbg}
\begin{longtable}{@{}L{4.8cm}L{2.2cm}L{1.6cm}L{4.8cm}@{}}
\toprule
\rowcolor{primary!15}
\textbf{Configuration} & \textbf{Readable?} & \textbf{HTTP} &
\textbf{A link the page emits} \\
\midrule
\endhead
setting off, rewrite off & no & 200 & \texttt{index.php?route=\dots} \\
\textbf{setting ON, rewrite off} & yes & \textbf{404} &
\texttt{/en-gb/catalog/desktops} \\
setting ON, rewrite ON & yes & 200 & \texttt{/en-gb/catalog/desktops} \\
\bottomrule
\end{longtable}

\smallskip
\textbf{The middle row is the finding.} Switching on the setting that every
guide tells you to switch on rewrites every link in the shop into a readable
one --- and \textbf{every one of them returns 404}. The home page still works,
because its own address is unchanged; everything reachable from it is gone.

\smallskip
\textbf{And the cause is \chref{security}'s.} The rewrite ships in
\texttt{.htaccess.txt}, and \texttt{AllowOverride None} means Apache never
reads \texttt{.htaccess} at all. One misconfiguration, found while looking for
an exposed error log, silently disables an unrelated feature two chapters
later.

\smallskip
\textbf{The application half is faultless.} Fetching
\texttt{index.php?\_route\_=htc-touch-hd} by hand returns 200 --- the lookup,
the keyword table and the routing all work. Only the server was never told.
\end{measured}

\begin{alertbox}[A change that makes a shop worse while appearing to improve
it]
Consider what the shopkeeper sees after ticking that box. The admin panel
reports success. The home page loads. The URLs in the address bar are now
pretty, which is visibly what was asked for. Nothing in any log is unusual,
because a 404 is a perfectly ordinary response.

\smallskip
What has actually happened is that \textbf{every product and category page in
the shop has stopped existing}, and the first person to find out will be a
customer or a crawler.

\smallskip
\textbf{The habit this argues for is one step long: after any configuration
change, follow a link.} Not the page you changed --- a link from it. The whole
of this failure is visible in one click and invisible in any amount of looking
at the setting.
\end{alertbox}

\section{Can a Crawler Tell the Pages Apart?}

\begin{measured}[Titles and descriptions across the catalogue]
From \texttt{code/m21\_seo.py --meta}, over all 19 product pages.

\smallskip
\begin{tabular}{@{}L{7.4cm}L{2.4cm}L{4.2cm}@{}}
Product pages checked & 19 & \\
Distinct titles & \textbf{19} & no duplicates --- good \\
Duplicate titles & 0 & \\
\textbf{Pages with no meta description} & \textbf{19} & \textbf{all of them} \\
Distinct descriptions & 0 & \\
\end{tabular}

\smallskip
\textbf{Titles are fine.} Each is the product's name, and because product names
differ, so do the titles. That is the one thing the platform gets right here,
and it gets it right by accident rather than by design --- a shop selling
``Blue T-Shirt'' in four sizes as four products would have four identical
titles.

\smallskip
\textbf{Not one page has a meta description.} The field exists in
\texttt{oc\_product\_description}; the demonstration data leaves it empty, and
the platform writes nothing in its place. A search engine will then compose a
description from whatever text it finds on the page, which for a product page
is usually the navigation.

\smallskip
\textbf{And no page carries a canonical link} --- not the home page, not a
category, not a product. So the same product reached through two categories is
two URLs with identical content, and nothing tells a crawler which one is the
real one.
\end{measured}

\begin{measured}[What a crawler is told, and how fast]
From \texttt{code/m21\_seo.py --crawl}.

\smallskip
\begin{tabular}{@{}L{6.4cm}L{3.0cm}L{4.4cm}@{}}
\texttt{robots.txt} & HTTP 200 & 16 \texttt{Disallow} rules \\
\dots names a sitemap? & \textbf{no} & \\
\texttt{/sitemap.xml} & \textbf{HTTP 404} & \\
A sitemap extension & \textbf{HTTP 404} & none ships \\
\end{tabular}

\smallskip
\begin{tabular}{@{}L{4.4cm}L{4.4cm}L{4.4cm}@{}}
& \textbf{Canonical?} & \textbf{First byte (median)} \\
Home & no & 114\,ms \\
Category & no & 133\,ms \\
Product & no & 110\,ms \\
\end{tabular}

\smallskip
\textbf{The \texttt{robots.txt} is good work.} Its sixteen rules block sorted,
paged and filtered variants of listing pages --- \texttt{?sort=},
\texttt{\&page=}, \texttt{?filter\_name=} --- which is exactly the faceted
navigation that otherwise generates thousands of near-identical URLs from
twenty products. Somebody thought about this.

\smallskip
\textbf{There is no sitemap of any kind}, and \texttt{robots.txt} does not name
one. A crawler therefore finds pages only by following links, so anything more
than a few hops from the home page is, in practice, not there. For a catalogue
of nineteen products this does not matter; at \chref{catalogue}'s fifty
thousand it is most of the shop.

\smallskip
\textbf{First byte is 110--133\,ms}, which is the server's own thinking time
and is respectable. \chref{ux} showed that the other 600\,KB is what the
customer actually waits for --- so this number, the one most often quoted, is
the smallest part of the experience.
\end{measured}

\section{The Order to Work In}

\begin{conceptbox}[Six things, each of which you can verify yourself]
\begin{tightnum}
  \item \textbf{Make the pages reachable.} Both halves of readable URLs, or
        neither. Then follow a link.
  \item \textbf{Publish a sitemap} and name it in \texttt{robots.txt}. The
        platform ships none, so this is a module (\chref{modules}) ---
        a query over products and categories emitting XML.
  \item \textbf{Give every page a canonical link.} One template override
        (\chref{ux}), and it ends the duplicate-content question for the whole
        shop.
  \item \textbf{Write meta descriptions}, starting with the twenty products
        that sell. It is data entry, it is nobody's favourite afternoon, and
        it is the only one of these six a search engine shows to a human.
  \item \textbf{Keep \texttt{robots.txt} blocking the facets.} It already does;
        the risk is somebody replacing it.
  \item \textbf{Then look at speed} --- and from \chref{ux}'s measurements,
        at page weight rather than at first byte.
\end{tightnum}

\smallskip
Everything above is checkable in your own shop, today, with no access to
anybody's analytics. Advice that cannot be checked that way goes after all six.
\end{conceptbox}

\begin{alertbox}[On the marketing half of this chapter's title]
The 2025 curriculum renames this course \emph{Digital Marketing \&
E-Commerce}, and a handout traced to the 2017 outline should say what it is not
covering. Campaign strategy, content marketing, social advertising and
attribution modelling are a different subject with different methods, and three
weeks of this course cannot do them justice.

\smallskip
What this course contributes to them is narrower and more durable:
\textbf{a shop that can be crawled, a coupon code per channel
(\chref{discounts}) so that a campaign is measurable at all, and the habit of
asking what the denominator is (\chref{dashboard}) before believing a rate.}
Without those three, a marketing budget is being spent on a shop that cannot
report what it bought.
\end{alertbox}

\begin{pitfall}
\begin{tight}
  \item \textbf{Switching on readable URLs without the rewrite.} Every link in
        the shop becomes a 404 and nothing reports it.
  \item \textbf{Changing a setting and checking the page you changed} instead
        of following a link from it.
  \item \textbf{Assuming \texttt{.htaccess} does anything on this stack.}
  \item \textbf{Leaving meta descriptions empty} and letting the engine write
        them from your navigation.
  \item \textbf{Shipping without canonical links}, so that one product at four
        URLs competes with itself.
  \item \textbf{Removing the \texttt{robots.txt} facet rules} to ``let
        everything be indexed''.
  \item \textbf{Optimising first byte} when \chref{ux} measured 600\,KB of
        assets after it.
  \item \textbf{Believing any claim about ranking that you cannot test.}
\end{tight}
\end{pitfall}

\begin{lab}[Make your shop findable, and break it first]
\begin{lstlisting}[language=bash]
# --- 1. Look at a link before changing anything ---------------------
cd code
curl -s "http://localhost:8090/index.php?route=product/category&path=20" \
  | grep -o 'href="http://localhost:8090/[^"]*"' | head -3

# --- 2. Turn on readable URLs the way every guide says --------------
# Admin: System > Settings > Edit > Server > Use SEO URLs = Yes.

# --- 3. Now FOLLOW A LINK -------------------------------------------
# Click any category. Write down what happened before reading on.

# --- 4. Measure all three configurations ----------------------------
python m21_seo.py --fix

# --- 5. Explain the middle row --------------------------------------
# Name the chapter that already measured the cause.

# --- 6. Audit what a crawler can tell apart -------------------------
python m21_seo.py --meta
python m21_seo.py --crawl

# --- 7. Fix the cheapest finding ------------------------------------
# Add a canonical link to the product template in your child theme
# (Chapter 8). One line. Re-run --crawl and confirm it appears.

# --- 8. Write the sitemap module ------------------------------------
# Chapter 19's structure: one controller, one query over enabled
# products and categories, XML out. Then name it in robots.txt.

# --- 9. Write ten meta descriptions ---------------------------------
# For ten real products. Then say honestly how long 50,000 would
# take, and what you would do instead.

# --- 10. Restore -----------------------------------------------------
python m21_seo.py --restore
\end{lstlisting}

\smallskip
\textbf{What to notice.} Step 3. Almost everybody ticks the box, sees a pretty
URL in the address bar, and concludes it worked. The entire shop is broken at
that moment, and one click shows it.
\end{lab}

\begin{checkpoint}
\begin{tightnum}
  \item A client turns on readable URLs and reports that ``the site looks the
        same but customers say pages are missing''. Explain precisely what has
        happened, name the two halves and which one is absent, give the command
        that demonstrates it, and say why nothing appeared in any log.
  \item All 19 product pages have distinct titles and none has a meta
        description. Say which of those two facts is luck and which is design,
        what happens on the pages with no description, and what you would do
        first for a catalogue of 50{,}000.
  \item This shop has no sitemap and no canonical links. Explain what each
        absence costs, say which matters more for a 19-product shop and which
        for a 50{,}000-product shop, and justify the difference.
\end{tightnum}
\end{checkpoint}

\begin{chaptersummary}
\begin{tight}
  \item Ask what is checkable: can a crawler read the pages, tell them apart,
        and find all of them? A shop failing that cannot rank whatever else is
        true.
  \item Readable URLs are two halves --- the application generating them and
        the server resolving them. Either alone is useless and the first alone
        is harmful.
  \item Measured: with the setting on and the rewrite off, every link the shop
        emits returns \textbf{404}, while the home page and
        \texttt{index.php?\_route\_=\dots} work perfectly. The cause is
        \chref{security}'s \texttt{AllowOverride None}.
  \item After any configuration change, follow a link. The whole of that
        failure is visible in one click and invisible in any amount of
        inspecting the setting.
  \item Measured: 19 distinct titles and 0 duplicates --- by luck, because
        product names differ --- and \textbf{19 pages with no meta description
        at all}, and no canonical link on any page.
  \item Measured: \texttt{robots.txt} blocks the faceted-navigation parameters,
        which is good work; there is no sitemap of any kind and nothing names
        one. First byte is 110--133\,ms, which \chref{ux} showed is the
        smallest part of what the customer waits for.
  \item Work in order: reachable, sitemap, canonical, descriptions, keep the
        facet rules, then speed --- and page weight before first byte.
  \item This course's contribution to marketing is a shop that can be crawled,
        a code per channel so campaigns are measurable, and always asking what
        the denominator is.
\end{tight}
\end{chaptersummary}

\begin{reviewq}
  \item \textbf{(MCQ)} With \texttt{config\_seo\_url} on and no server rewrite,
        a link the category page emits returns: (a)~200; (b)~301; (c)~404;
        (d)~500.
  \item \textbf{(MCQ)} How many of the 19 product pages carried a meta
        description? (a)~all 19; (b)~about half; (c)~one; (d)~none.
  \item \textbf{(Short)} Explain what a canonical link is for, with an example
        of two URLs in this shop that need one.
  \item \textbf{(Short)} Why does \texttt{robots.txt} block \texttt{?sort=} and
        \texttt{?filter\_name=}, and what would happen if it did not?
  \item \textbf{(Short)} Give the one-step habit this chapter argues for after
        any configuration change, and the failure it would have caught.
  \item \textbf{(Applied)} You inherit a shop with 8{,}000 products, no
        sitemap, no canonical links and meta descriptions on 200 of them. Set
        out your plan in order, with the effort and the expected benefit of
        each step, and say which steps you would automate and how.
  \item \textbf{(Applied)} Specify the sitemap module this platform lacks: what
        it queries, how it handles a 50{,}000-product catalogue within the
        sitemap size limits, how often it is regenerated, how it is named in
        \texttt{robots.txt}, and how you would verify it works without waiting
        for a search engine.
\end{reviewq}
